%%%%%%%%%%%%
% PACKAGES %
%%%%%%%%%%%%

\usepackage[utf8]{inputenc}
\usepackage[T1]{fontenc}
\usepackage{textcomp}
\usepackage{csquotes}
\usepackage{url}
\usepackage{graphicx}
\usepackage{float}
\usepackage{booktabs}
\usepackage{enumitem}
\usepackage{import}
\usepackage{xifthen}
\usepackage{pdfpages}
\usepackage{transparent}
\usepackage{parskip}
\usepackage{emptypage}
\usepackage{subcaption}
\usepackage{multicol}
\usepackage{cmbright}
\usepackage{xcolor}
\usepackage{amsmath}
\usepackage{amsfonts}
\usepackage{mathtools}
\usepackage{amsthm}
\usepackage{amssymb}
\usepackage{mathrsfs}
\usepackage{cancel}
\usepackage{bm}
\usepackage{systeme}
\usepackage{stmaryrd}
\usepackage{siunitx}
\usepackage{mdframed}
\usepackage{etoolbox}
\usepackage{fancyhdr}
\usepackage{todonotes}
\usepackage{tcolorbox}
\usepackage{fourier, erewhon, cabin}
\usepackage[english]{babel}
\usepackage{titletoc}
\usepackage{lipsum}
\usepackage[utf8]{inputenc}
\usepackage[T1]{fontenc}
\usepackage{textcomp}
\usepackage[nottoc]{tocbibind}
\usepackage{tikz}
\usepackage{tikz-cd}
\usepackage{thmtools}
\usepackage{thm-restate}
\usepackage{multirow}
\usepackage{hyperref}
\usepackage[noabbrev]{cleveref}
\usepackage{ulem}
\usepackage[
    sorting=nyt,
    style=alphabetic
]{biblatex}


%%%%%%%%%%%%%%%%%%%%%%%%%%%%%%%%%%%%
% New commands and Re-New commands %
%%%%%%%%%%%%%%%%%%%%%%%%%%%%%%%%%%%%

\newcommand{\newtodo}[1]{\todo[color=orange]{TODO:                              \hspace{26.5pt}#1.}}
\newcommand{\canwait}[1]{\todo[color=green]{CAN WAIT:                           \hspace{9.5pt}#1.}}
\newcommand{\important}[1]{\todo[color=yellow]{IMPORTANT:                       #1.}}
\newcommand{\urgent}[1]{\todo[color=red]{URGENT:                                \hspace{16.5pt}#1.}}
\newcommand{\done}[1]{\todo[disable]{#1}\addcontentsline{tdo}{todo}{\sout{DONE: \hspace{26.5pt}#1.}}}

\newcommand\contra{\scalebox{1.5}{$\lightning$}}
\newcommand\correct[2]{\ensuremath{\:}{\color{red}{#1}}\ensuremath{\to }{\color{correct}{#2}}\ensuremath{\:}}
\newcommand\green[1]{{\color{correct}{#1}}}
\newcommand\hr{
    \noindent\rule[0.5ex]{\linewidth}{0.5pt}                        % Horizontal rule
}
\newcommand\hide[1]{}                                               % Hide parts
\setcounter{secnumdepth}{0}

\def\@lesson{}
\def\@lessonabbr{}
\newcommand{\lesson}[3]{
    \ifthenelse{\isempty{#3}}{
        \def\@lesson{Lesson #1}
        \subsection[\@lesson]{\hr \@lesson \hfill \small{#2}}
    }{
        \def\@lesson{Lesson #1: #3}
        \subsection[\@lesson]{\hr \@lesson \hfill \small{#2}}
    }
    \def\@lessonabbr{Les #1}
}

\def\@lecture{}
\def\@lectureabbr{}
\newcommand{\lecture}[3]{
    \ifthenelse{\isempty{#3}}{
        \def\@lecture{Lecture #1}
        \subsection[\@lecture]{\hr \@lecture \hfill \small{#2}}
    }{
        \def\@lecture{Lecture #1: #3}
        \subsection[\@lecture]{\hr \@lecture \hfill \small{#2}}
    }
    \def\@lectureabbr{Lec #1}
}

\newcommand\circled[1]{
    \begin{tikzpicture}[baseline=(char.base)]%
        \node[circle,draw,inner sep=1pt] (char) {\textsf{#1}};%
\end{tikzpicture}}
\newcommand\mc[1]{\footnotesize\circled{#1}}
\newcommand{\incfig}[1]{%
    \def\svgwidth{\columnwidth}
    \import{./figures/}{#1.pdf_tex}
}

\let\oldcite\cite
\renewcommand*\cite[1]{\textsubscript{ (\textit{\oldcite{#1}}) }}

\newcommand{\C}{\ensuremath{\mathbb{C}}}
\newcommand{\Q}{\ensuremath{\mathbb{Q}}}
\renewcommand{\O}{\ensuremath{\mathbb{O}}}
\newcommand{\Z}{\ensuremath{\mathbb{Z}}}
\newcommand{\R}{\ensuremath{\mathbb{R}}}
\newcommand{\N}{\ensuremath{\mathbb{N}}}

\let\implies\Rightarrow
\let\impliedby\Leftarrow
\let\iff\Leftrightarrow
\let\epsilon\varepsilon

\newcommand{\fullwidthincfig}[2][0.90]{%
    % \ifworking{\makebox[0pt][l]{\color{gray}{\scriptsize\textsf{#2}}}}\fi%
    \def\svgwidth{#1\paperwidth}
    \import{./figures/}{#2.pdf_tex}
}
\newcommand{\minifig}[2]{%
    \def\svgwidth{#1}%
    \begingroup%
    \setbox0=\hbox{\import{./figures/}{#2.pdf_tex}}%
    \parbox{\wd0}{\box0}\endgroup%
    \hspace*{0.2cm}
}
\newcommand\charlotte[1]{{\color{blue}{#1}}}

\def\block(#1,#2)#3{\multicolumn{#2}{c}{\multirow{#1}{*}{$ #3 $}}}


%%%%%%%%%%%%
% Settings %
%%%%%%%%%%%%

\sisetup{locale = FR}                                               % Si unitx
\def\thm@space@setup{ \thm@preskip=\parskip \thm@postskip=0pt }     % Fix some spacing
\let\svlim\lim\def\lim{\svlim\limits}
\tcbuselibrary{breakable}                                           % Make boxes breakable
\pdfsuppresswarningpagegroup=1                                      % Fix some stuff
\addbibresource{bibliography.bib}
\hypersetup{                                                        % Make table of contents clickable
    colorlinks,
    linkcolor={black},
    citecolor={black},
    urlcolor={blue!80!black}
}
\mdfsetup{skipabove=1em,skipbelow=0em, innertopmargin=12pt, innerbottommargin=8pt}
\renewcommand{\sectionmark}[1]{\markboth{#1}{}}                     % Set the \leftmark
\fancypagestyle{plain}{                                             % Changes footer and header
    \fancyhf{}

    \fancyhead[R,L]{\hr}
    \fancyhead[C]{\leftmark}                                        % Center
    \fancyhead[RO,LE]{\@lessonabbr}                                 % Right odd,  Left even
    \fancyhead[RE,LO]{\thepage}                                     % Right even, Left odd
    
    \fancyfoot[RE,LO]{\theTitle}
    \fancyfoot[RO,LE]{\thedate}
    \fancyfoot[C]{\theauthor}

    \renewcommand{\headrulewidth}{0.4pt}
    \renewcommand{\footrulewidth}{0.4pt}
}


%%%%%%%%%%%%%%%%%%%%
% New environments %
%%%%%%%%%%%%%%%%%%%%

\mdfsetup{skipabove=1em,skipbelow=0em}
\theoremstyle{definition}
\newmdtheoremenv{lemma}{Lemma}
\newmdtheoremenv{prop}{Proposition}
\newmdtheoremenv{theorem}{Theorem}
\newmdtheoremenv{postulate}{Postulate}
\newmdtheoremenv{myproof}{Proof}
\newmdtheoremenv{definition}{Definition}
\newmdtheoremenv{note}{Note}
\newmdtheoremenv{problem}{Problem}
\newmdtheoremenv{question}{Question}
\newmdtheoremenv{corollary}{Corollary}
\newmdtheoremenv{example}{Example}
\newmdtheoremenv{notation}{Notation}
\newmdtheoremenv{repetition}{Repetition}
\newmdtheoremenv{property}{Property}
\newmdtheoremenv{intuition}{Intuition}
\newmdtheoremenv{observation}{Observation}
\newmdtheoremenv{conclusion}{Conclusion}
\declaretheorem[numbered=no, style=todo, name=TODO]{TODO}

\newenvironment{remark}{\begin{tcolorbox}[
    arc=0mm,
    colback=white,
    colframe=purple,
    title=Remark,
    fonttitle=\sffamily,
    breakable
]}{\end{tcolorbox}}
\newenvironment{answer}{\begin{tcolorbox}[
    arc=0mm,
    colback=white,
    colframe=blue,
    title=Answer,
    fonttitle=\sffamily,
    breakable
]}{\end{tcolorbox}}
\newenvironment{info}{\begin{tcolorbox}[
    arc=0mm,
    colback=white,
    colframe=gray,
    title=Info,
    fonttitle=\sffamily,
    breakable
]}{\end{tcolorbox}}
\newenvironment{terminology}{\begin{tcolorbox}[
    arc=0mm,
    colback=white,
    colframe=green!60!black,
    title=Terminology,
    fonttitle=\sffamily,
    breakable
]}{\end{tcolorbox}}
\renewenvironment{warning}{\begin{tcolorbox}[
    arc=0mm,
    colback=white,
    colframe=red,
    title=Warning,
    fonttitle=\sffamily,
    breakable
]}{\end{tcolorbox}}
\newenvironment{caution}{\begin{tcolorbox}[
    arc=0mm,
    colback=white,
    colframe=yellow,
    title=Caution,
    fonttitle=\sffamily,
    breakable
]}{\end{tcolorbox}}

\definecolor{correct}{HTML}{009900}
\theoremstyle{definition}
\declaretheoremstyle[headfont=\bfseries\sffamily, bodyfont=\normalfont, mdframed={ nobreak } ]{thmgreenbox}
\declaretheoremstyle[headfont=\bfseries\sffamily, bodyfont=\normalfont, mdframed={ nobreak } ]{thmredbox}
\declaretheoremstyle[headfont=\bfseries\sffamily, bodyfont=\normalfont]{thmbluebox}
\declaretheoremstyle[headfont=\bfseries\sffamily, bodyfont=\normalfont]{thmblueline}
\declaretheoremstyle[headfont=\bfseries\sffamily, bodyfont=\normalfont, numbered=no, mdframed={ rightline=false, topline=false, bottomline=false, }, qed=\qedsymbol ]{thmproofbox}
\declaretheoremstyle[headfont=\bfseries\sffamily, bodyfont=\normalfont, numbered=no, mdframed={ nobreak, rightline=false, topline=false, bottomline=false } ]{thmexplanationbox}
\declaretheoremstyle[headfont=\bfseries\sffamily, bodyfont=\normalfont, numbered=no, mdframed={ nobreak, rightline=false, topline=false, bottomline=false } ]{thmexplanationbox}
\declaretheorem[sibling=definition, style=thmbluebox,  name=Nonexample]{noneg}
\declaretheorem[numbered=no, style=thmexplanationbox, name=Proof]{explanation}
\declaretheorem[numbered=no, style=thmproofbox, name=Proof]{replacementproof}
\declaretheorem[style=thmbluebox,  numbered=no, name=Exercise]{ex}
\declaretheoremstyle[
    headfont=\bfseries\sffamily\color{RawSienna!70!black}, bodyfont=\normalfont,
    mdframed={
        linewidth=2pt,
        rightline=false, topline=false, bottomline=false,
        linecolor=RawSienna, backgroundcolor=RawSienna!5,
    }
]{todo}
\AtBeginEnvironment{quote}{\itshape}
